\chapter{Trabalhos Relacionados}
\label{cap:trabalhos-relacionados}

Este capítulo apresenta e discute pesquisas que analisam como a integração de diferentes
áreas do conhecimento contribui para a eficácia no desenvolvimento de soluções e para a
elevação da qualidade do aprendizado. O objetivo desta revisão é demonstrar que a
superação da fragmentação curricular, por meio da articulação entre saberes técnicos,
científicos e sociais, constitui um caminho estratégico para a resolução de problemas
complexos da sociedade contemporânea.

\section{A interdisciplinaridade como estratégia para a formação integral no ensino
superior - Ribeiro et al., 2025.
}\label{sec:ribeiro-2025}

O trabalho de Ribeiros e colegas (2025) discute como a união de diferentes áreas do
conhecimento ajuda na formação completa dos alunos na universidade. Os autores
explicam que, hoje em dia, o ensino ainda é muito dividido em "caixinhas" (disciplinas
isoladas), o que faz com que os estudantes tenham dificuldade de resolver problemas
reais, que são sempre complexos e misturados. Para mudar isso, o artigo defende o uso de
metodologias ativas, como a Aprendizagem Baseada em Problemas (PBL), onde o aluno
deixa de ser um espectador e passa a ser o centro do aprendizado.

A pesquisa mostra que, quando as matérias se conversam e o aluno tem um papel ativo,
ele desenvolve melhor o pensamento crítico e a capacidade de criar soluções que
realmente funcionam. Os resultados indicam que essa integração não serve apenas para
passar de ano, mas para preparar o profissional para os desafios de um mercado que exige
visão de mundo e flexibilidade.

Este estudo é uma revisão bibliográfica de natureza integrativa que analisou publicações e
documentos produzidos entre 2015 e 2025. Como resultado principal, a análise teórica
consolidou o consenso de que abordagens estritamente disciplinares são insuficientes
para a resolução de problemas complexos. O trabalho demonstra, por meio da literatura
examinada, que currículos integrados e orientados a situações reais fortalecem o
engajamento estudantil e a articulação entre ensino, pesquisa e extensão, validando o
protagonismo do aluno como o elo fundamental entre a teoria acadêmica e a entrega de
soluções assertivas na prática profissional.

\subsection{Integração de Inovação Social, Empreendedorismo, Propriedade Intelectual e
Desenvolvimento Científico: Uma Abordagem Multidisciplinar - Bentes, 2024.} \label{sec:bentes-2024}

Neste estudo, Bentes (2024) analisa a convergência entre inovação social,
empreendedorismo, propriedade intelectual e desenvolvimento científico no ambiente
universitário. O autor identifica que a fragmentação dos currículos e o isolamento entre as
disciplinas impedem que as universidades respondam de forma prática às necessidades
do mercado e da sociedade. Para fundamentar essa análise, foi realizada uma revisão
integrativa da literatura de 1996 a 2023, utilizando o papel dos Núcleos de Inovação
Tecnológica (NITs) como ponto central para observar como a gestão da propriedade
intelectual pode servir de ponte entre a pesquisa acadêmica e a criação de produtos e
serviços reais.

Os resultados indicam que a integração dessas quatro áreas permite que as instituições
de ensino superior atuem como facilitadoras de desenvolvimento econômico e social. O
estudo aponta que a gestão estratégica feita pelos NITs garante a proteção das invenções
e facilita a transferência de tecnologia para o mercado, o que viabiliza a aplicação de
soluções para problemas sociais complexos. Dessa forma, a multidisciplinaridade é
apresentada como a estratégia necessária para formar profissionais capazes de conectar
o conhecimento científico a soluções inovadoras e sustentáveis. 

Esse trabalho apresenta pontos de convergência com a presente monografia, pois ambos
defendem que a integração de diferentes saberes é o melhor caminho para preparar o
estudante para desafios reais. A diferença é que, enquanto Bentes foca mais em como a
universidade deve se organizar como instituição, esta pesquisa vai além e foca no
protagonismo do aluno, defendendo que o estudante ser o centro do próprio aprendizado
é o que garante que ele realmente aprenda com qualidade e desenvolva soluções mais
assertivas.

\subsection{Desenvolvimento de Software para Solução Numérica de Equações Polinomiais:
Uma Abordagem Interdisciplinar – Silva et. al., 2018.} \label{sec:silva-2018}

O estudo de Silva e colaboradores (2018) investiga como a multidisciplinaridade impacta o
processo de design de interfaces digitais. Os autores partem do problema de que o design
moderno exige a colaboração de profissionais de diferentes áreas (como computação,
artes e psicologia), mas que essa interação frequentemente falha devido à falta de uma
linguagem comum e processos integrados. Essa fragmentação gera soluções que não
atendem plenamente às necessidades do usuário final.

Para analisar esse cenário, os autores realizaram um estudo de caso prático envolvendo
equipes multidisciplinares no desenvolvimento de produtos digitais. A análise foi feita por
meio de observação direta e entrevistas com os envolvientes, buscando entender como a
troca de conhecimentos influenciava a qualidade do artefato final. Os resultados mostram
que a integração de saberes não ocorre de forma espontânea, ela exige métodos que
facilitem o diálogo entre as áreas. Quando essa barreira é superada, observa-se um
aumento significativo na assertividade do design, resultando em interfaces mais
acessíveis e funcionais.

O trabalho de Silva e colegas (2018) é fundamental para esta monografia por apresentar
evidências práticas de que a multidisciplinaridade é um desafio de comunicação e
processo. Ele reforça a tese de que unir diferentes visões melhora a qualidade do software.
Esta pesquisa pretende avançar ao demonstrar que o protagonismo do aluno, quando
inserido nesse contexto de integração, é o que potencializa a aprendizagem e garante a
criação de soluções tecnológicas superiores.


\begin{table}[h!]	
	\centering
	\Caption{\label{tab:comparacao-trabalho-relacionados} Comparação dos trabalhos relacionados.}	
	\UFCtab{}{
		\begin{tabular}{p{2.8cm}p{3cm}p{3cm}p{3cm}p{3cm}}
			\toprule
			\textbf{Critério} & \textbf{Ribeiro et al. (2025)} & \textbf{Bentes (2024)} & \textbf{Silva et al. (2018)} & \textbf{Esta Monografia} \\
			\midrule \midrule
			\textbf{Foco Central} & Estratégias para formação integral no ensino superior. & Integração de inovação, gestão e ciência. & Processo de design e colaboração entre áreas. & Análise de como a integração de conteúdo possibilita o protagonismo.\\
			\textbf{Objeto de estudo} & Metodologias ativas e interdisciplinaridade. & NITs e currículos de Ensino Superior. & Equipes de design e desenvolvimento digital. & Metodologias ativas e integração multidisciplinar.\\
			\textbf{Abordagem metodológica} & Revisão teórica e análise de estratégias. & Revisão integrativa de literatura. & Estudo de caso com observação e entrevistas. & Análise de estudo de caso: ArenaHub.\\
			\textbf{Integração como motor de solução} & Sim (foco no desenvolvimento humano). & Sim (soluções para desafios sociais). & Sim (criação de interfaces mais assertivas). & Sim (soluções tecnológicas e educacionais).\\
			\textbf{Protagonismo do Aluno/Sujeito} & Sim (fator chave no PBL). & Indireto (foco na formação profissional). & Parcial (foco na equipe profissional). & Sim (fator central da qualidade do ensino).\\
			\bottomrule
		\end{tabular}
	}{
	\Fonte{Elaborado pelo autor.}
} 
\end{table}